\documentclass[11pt,a4paper]{article}
\usepackage[utf8]{inputenc}
\usepackage[T1]{fontenc}
\usepackage[spanish]{babel}
\usepackage{lmodern}
\usepackage[margin=2.2cm]{geometry}
\usepackage{amsmath,amssymb}
\usepackage{booktabs}
\usepackage{xcolor}
\usepackage{enumitem}
\usepackage{fancyhdr}

\setlength{\parindent}{0pt}
\setlength{\parskip}{0.55em}
\pagestyle{fancy}
\fancyhf{}
\lhead{Ejemplo de SuperMinHash}
\rhead{Otmar Ertl, 2017}
\cfoot{\thepage}

\title{SuperMinHash paso a paso\\\large Ejemplo de construcción de una firma}
\author{}
\date{}

\begin{document}
\maketitle
\vspace{-2.0em}

\section*{1. Qué vamos a calcular}

SuperMinHash resume un conjunto de elementos en una \textbf{firma} de tamaño
$m$. Si $m=3$, la firma tiene tres posiciones:
\[
  h = [h_0,h_1,h_2].
\]
Cada posición conserva el menor valor que se le haya asignado. En una iteración
con índice $j$, se genera un candidato $r+j$, donde $r$ es un número aleatorio
entre 0 y 1. Una permutación aleatoria determina a qué posición de la firma se
asigna ese candidato.

\textbf{Importante:} los valores que siguen son una ejecución inventada para
explicar el mecanismo. No son una salida fija del algoritmo.

\section*{2. Procesamiento del primer elemento}

La firma parte sin valores finitos: $h=[\infty,\infty,\infty]$. Para simplificar,
mostramos directamente qué posición selecciona la permutación en cada paso.

\begin{center}
\begin{tabular}{ccccccl}
\toprule
Iteración $j$ & Aleatorio $r$ & Candidato $r+j$ & Posición & Valor previo & Valor nuevo & Acción \\
\midrule
0 & 0,4 & 0,4 & $h_2$ & $\infty$ & 0,4 & Actualiza \\
1 & 0,2 & 1,2 & $h_0$ & $\infty$ & 1,2 & Actualiza \\
2 & 0,7 & 2,7 & $h_1$ & $\infty$ & 2,7 & Actualiza \\
\bottomrule
\end{tabular}
\end{center}

La firma queda $h=[1,2;\,2,7;\,0,4]$. Se hicieron las tres iteraciones
posibles. La parte entera de sus valores es, respectivamente, 1, 2 y 0.

\section*{3. El histograma permite cortar el segundo elemento}

El algoritmo mantiene un histograma con la cantidad de posiciones de la firma
que tienen cada parte entera. La variable $a$ guarda la parte entera más alta
que todavía aparece. Después del primer elemento:

\begin{center}
\begin{tabular}{cccc}
\toprule
Firma actual & Conteo de parte entera 0 & Conteo de 1 & Conteo de 2 \\
\midrule
$[1,2;\,2,7;\,0,4]$ & 1 & 1 & 1 \\
\bottomrule
\end{tabular}
\qquad $a=2$
\end{center}

Para el segundo elemento:

\begin{enumerate}[leftmargin=1.5em]
  \item En $j=0$, el candidato es 0,1 y se asigna a $h_1$. Como
        $0,1<2,7$, actualiza la firma a $[1,2;\,0,1;\,0,4]$.
  \item La parte entera 2 ya no aparece: las partes enteras son ahora
        $[1,0,0]$. El histograma lo refleja y baja $a$ de 2 a 1.
  \item En $j=1$, el candidato es 1,6 y se asigna a $h_2$. No actualiza,
        porque $1,6>0,4$.
  \item La siguiente iteración sería $j=2$. Como $j>a$ (2>1), el algoritmo
        se detiene para este elemento.
\end{enumerate}

\textbf{¿Por qué puede parar?} Cuando $j=2$, cualquier candidato es al menos
2. Todas las posiciones actuales son menores que 2, así que ningún candidato
restante podría reemplazar un mínimo por uno menor.

\section*{4. Un tercer elemento muestra un corte más temprano}

La firma actual es $[1,2;\,0,1;\,0,4]$, con $a=1$. Para el tercer elemento,
supongamos que en $j=0$ se genera el candidato 0,2 y se asigna a $h_0$. Como
$0,2<1,2$, la firma pasa a $[0,2;\,0,1;\,0,4]$.

Ahora las tres partes enteras son 0, de modo que $a=0$. La iteración siguiente
sería $j=1$, pero $j>a$; el algoritmo omite $j=1$ y $j=2$ para este elemento.

\section*{5. ¿Dónde está el ahorro?}

\begin{center}
\begin{tabular}{lcc}
\toprule
Elemento procesado & Iteraciones realizadas & Iteraciones omitidas de un máximo de 3 \\
\midrule
Primero & 3 & 0 \\
Segundo & 2 & 1 \\
Tercero & 1 & 2 \\
\midrule
Total ilustrativo & \textbf{6} & \textbf{3} \\
\bottomrule
\end{tabular}
\end{center}

Un recorrido completo habría realizado 9 iteraciones en este ejemplo; el corte
temprano deja 6. Al omitir iteraciones, se evitan generar números aleatorios,
hacer intercambios de la permutación y probar candidatos que no pueden mejorar
la firma. El ahorro depende de los valores aleatorios y de los datos; esta
ejecución pequeña solo ilustra el mecanismo.

El histograma se actualiza cuando cambia una posición. Por eso, el algoritmo no
necesita recorrer de nuevo todas las posiciones para encontrar $a$ en cada
iteración. Sí hay un costo inicial para crear las $m$ posiciones y las
estructuras auxiliares.

\section*{6. Qué no significa “mejorar”}

\begin{itemize}[leftmargin=1.5em]
  \item “Mejorar la firma” significa únicamente bajar el mínimo guardado en
        alguna posición; no significa que dos documentos se parezcan más.
  \item El algoritmo no compara primero las partes enteras de dos firmas.
        Para estimar Jaccard se comparan los valores completos, posición por
        posición.
  \item Las posiciones de la firma no son los \emph{buckets} de LSH
        (\emph{Locality-Sensitive Hashing}), que se usan después para agrupar
        candidatos.
\end{itemize}

\textbf{Idea central:} SuperMinHash mantiene los mínimos de la firma y usa un
histograma incremental para saber cuándo los candidatos restantes ya no
pueden reducir ninguno de ellos.

\end{document}
